%% ma: L11-B15-0063
%% lop: 11
%% chuong: 5
%% bai: 15
%% dang: 11.15.9
%% loai: TN4
%% muc: VD
%% dap_an: "C"
%% nguon: "tai-lieu-hoc-tap-mon-toan-lop-11-theo-chuong-trinh-giao-duc-pho-thong-2018.pdf (Lớp Toán Thầy Duy, Phạm Lê Duy, THPT Chu Văn An), Chương 3 Bài 1, Phần 1, Câu 52 (tr.332)"
%% kien_thuc: "Lớp 11 Bài 15 (giới hạn của dãy số)"
%% trang_thai: chinh-thuc
%% ghi_chu: "Đáp án gốc A sai (khi n≥3 chiều cao 175 cm vượt 1,5 m); đúng là C. Đã duyệt 03/11/2026."
\begin{ex}
Một mô hình gồm các khối cầu xếp chồng lên nhau tạo thành một cột thẳng đứng. Biết rằng mỗi khối cầu có bán kính gấp đôi khối cầu nằm ngay trên nó và bán kính khối cầu dưới cùng là $50$ cm. Hỏi mệnh đề nào sau đây là \textbf{đúng}?
\choice{Chiều cao mô hình không quá $1{,}5$ mét.}{Chiều cao mô hình tối đa là $2$ mét.}{\True Chiều cao mô hình dưới $2$ mét.}{Mô hình có thể đạt được chiều cao tùy ý.}
\loigiai{\sloppy Đường kính các khối cầu từ dưới lên: $100;50;25;\dots$ (cấp số nhân công bội $\dfrac12$). Chiều cao của $n$ khối cầu là \[100\cdot\dfrac{1-\left(\frac12\right)^n}{1-\frac12}=200\left[1-\left(\tfrac12\right)^n\right]<200\] (cm), và tiến dần tới $200$ cm $=2$ m nhưng không bao giờ đạt. Vậy chiều cao luôn dưới $2$ m; khi $n\ge3$ chiều cao đã vượt $1{,}5$ m nên A sai; B sai vì không đạt đúng $2$ m.}
\end{ex}
